\chapter{Streams}
\label{ch:streams}

\standfirst{A stream is a position in a collection. Reading and writing become
the same two verbs regardless of what is underneath, which is why building a
string, parsing input and talking to a file all look alike.}

\section{Why streams}

Building a string by repeated concatenation copies everything each time:

\begin{st}
"quadratic, and allocates a new String on every iteration"
| s |
s := ''.
#('a' 'b' 'c') do: [:each | s := s , each].
\end{st}

A stream keeps a position and grows its target as needed:

\begin{st}
| s |
s := WriteStream on: (String new: 0).
#('a' 'b' 'c') do: [:each | s nextPutAll: each].
s contents                        "'abc'"
\end{st}

That is the idiom you will write hundreds of times. \ct{printOn:} methods,
report generation, and every bit of serialisation in the library are this
shape.

\section{The three kinds}

\begin{center}
\small
\begin{tabular}{@{}lp{0.60\textwidth}@{}}
\toprule
\ctp{ReadStream} & reads from an existing collection\\
\ctp{WriteStream} & writes into a collection, growing it\\
\ctp{ReadWriteStream} & both, with one position\\
\bottomrule
\end{tabular}
\end{center}

\begin{st}
ReadStream on: #(1 2 3)
ReadStream on: 'hello'
WriteStream on: (String new: 0)
WriteStream on: (Array new: 0)
WriteStream with: 'existing'       "positioned at the end, appends"
ReadWriteStream on: (String new: 0)
\end{st}

\begin{stwarn}
\ct{WriteStream on: aCollection} discards what is already there and writes
from position zero. \ct{WriteStream with: aCollection} keeps it and appends.
The names give no hint which is which, and choosing wrong silently loses your
data.
\end{stwarn}

\section{Reading}

\begin{st}
| s |
s := ReadStream on: #(1 2 3 4).
s next                     "1"
s peek                     "2 -- look without consuming"
s skip: 1
s next                     "3"
s atEnd                    "false"
s upToEnd                  "(4 )"
s reset                    "back to the beginning"
s next: 2                  "(1 2 ) -- the next two"
s upTo: 3                  "everything before the first 3, consuming it"
s do: [:each | ... ]       "the rest, one at a time"
\end{st}

\ct{next} past the end answers \ct{nil} rather than raising, so
\ct{[s atEnd] whileFalse: [...]} is the loop to write.

\section{Writing}

\begin{st}
| s |
s := WriteStream on: (String new: 0).
s nextPut: $a.                 "one element"
s nextPutAll: 'bcd'.           "many"
s print: 42.                   "the printString of anything"
s nextPutAll: 42 printString.  "the same thing, written out"
s space; tab; cr.              "convenience, on character streams"
s contents                     "everything written so far"
\end{st}

\ct{print:} is the one that saves the most typing, and it works for any object
because it goes through \ct{printOn:}. If you have given your class a
\ct{printOn:}, it prints properly inside collections, inside other objects'
\ct{printOn:} methods, and here.

\section{Streams and \texttt{printOn:}}

This is the connection worth seeing clearly. \ct{printString} is defined
roughly as:

\begin{st}
printString
    | s |
    s := WriteStream on: (String new: 0).
    self printOn: s.
    ^s contents
\end{st}

So writing \ct{printOn:} for your class gets you \ct{printString},
\ct{displayString}, correct printing inside collections, and useful output in
the Inspector and Debugger---all from one method that never builds a string.

\begin{st}
printOn: aStream
    aStream nextPutAll: 'an Account of '; print: balance
\end{st}

\section{What is not here}

\ct{collect:} is the one piece of collection protocol a stream has, and it
is eager: it drains the receiver and answers a \ct{ReadStream} on an
\ct{Array} of the results.

\begin{st}
((ReadStream on: #(1 2 3)) collect: [:x | x * 2]) upToEnd   "#(2 4 6)"
\end{st}

\ct{select:}, \ct{reject:}, \ct{detect:}, \ct{inject:into:} and
\ct{includes:} are not there at all, and \ct{size} on a \ct{ReadStream}
answers zero rather than the size of what it reads---ask its \ct{contents}
for that. Read into a collection first, or use \ct{do:}.

\section{Files}

Files are streams too. \ct{Disk} is the global for the file system, and
\ct{FileStream} is what you get:

\begin{st}
| f |
f := FileStream fileNamed: 'notes.txt'.
f nextPutAll: 'hello'; close.

(FileStream fileNamed: 'notes.txt') contentsOfEntireFile
\end{st}

The File List window---\ct{file list} on the desktop menu---is a browser for
the same thing, and Chapter~\ref{ch:tools} covers it.

Underneath, this is a small set of POSIX primitives the VM supplies: open,
close, read, write, truncate, size, remove, rename and list. A file's
descriptor lives in an instance variable, which means it is part of the object
memory and gets written into a snapshot; the VM only believes a descriptor
that the current process opened, so a resumed image reopens by name rather
than writing to whatever that number now means.

\section{The \texttt{Transcript}}

\ct{Transcript} is a stream you can write to from anywhere, and it appears in
the System Transcript window if one is open.

\begin{st}
Transcript show: 'hello'; cr.
Transcript show: 42 printString; cr.
Transcript display: 42; cr.
\end{st}

\ct{show:} wants a String. \ct{display:} takes anything and uses its
\ct{displayString}, which is usually what you meant.

\begin{stwarn}
Output from concurrent processes may interleave. If you want a line to arrive
whole, hold the Transcript's monitor across the entry, or build the line in a
\ct{WriteStream} and \ct{show:} it in one send. See
Chapter~\ref{ch:contract}.
\end{stwarn}
